\documentclass[a4paper,12pt]{article}
\usepackage[utf8]{inputenc}
\usepackage[T2A]{fontenc}
\usepackage[russian]{babel}
\usepackage{geometry}
\usepackage{graphicx}
\usepackage{amsmath}
\usepackage{amsfonts}
\usepackage{hyperref}
\usepackage{array}
\usepackage{longtable}
\usepackage{booktabs}
\usepackage{caption}
\usepackage{listings}
\usepackage{xcolor}
\usepackage{float}

\geometry{top=2cm, bottom=2cm, left=2.5cm, right=1.5cm}

\definecolor{codegreen}{rgb}{0,0.6,0}
\definecolor{codegray}{rgb}{0.5,0.5,0.5}
\definecolor{codeblue}{rgb}{0,0,0.95}
\definecolor{backcolour}{rgb}{0.95,0.95,0.92}

\lstdefinestyle{mystyle}{
    backgroundcolor=\color{backcolour},
    commentstyle=\color{codegreen},
    keywordstyle=\color{codeblue},
    numberstyle=\tiny\color{codegray},
    stringstyle=\color{codegreen},
    basicstyle=\ttfamily\footnotesize,
    breakatwhitespace=false,
    breaklines=true,
    captionpos=b,
    keepspaces=true,
    numbers=left,
    numbersep=5pt,
    showspaces=false,
    showstringspaces=false,
    showtabs=false,
    tabsize=2
}

\lstset{style=mystyle}

\title{Практическая работа 2 по курсу\\
"Базы данных"\\
Вариант 7\\
Модель "Заказ билетов"}
\author{Мурадян Артём Андреевич\\
БПМ232}
\date{\today}

\begin{document}

\maketitle
\newpage

\tableofcontents
\newpage

\section{Описание задания}

В данной практической работе необходимо выполнить следующие задачи:
\begin{enumerate}
    \item Сформировать запросы к базе данных в виде представлений и хранимых процедур с демонстрацией всех возможных результатов (пустая таблица, одна строка, несколько строк).
    \item Провести оптимизацию одного из запросов и продемонстрировать результат.
    \item Создать триггер для справочника с сохранением истории изменений.
    \item Создать пользователей и разграничить права доступа.
    \item Создать представления для каждой группы пользователей.
    \item Создать резервную копию БД, удалить и восстановить её.
\end{enumerate}

\section{Запросы в виде представлений и хранимых процедур}

\subsection{Представление: Список пассажиров заданного рейса}

\begin{lstlisting}[language=SQL, caption=Создание представления v\_passengers\_by\_flight]
CREATE OR REPLACE VIEW v_passengers_by_flight AS
SELECT 
    f.flight_number,
    f.departure_city,
    f.arrival_city,
    f.departure_time,
    p.last_name,
    p.first_name,
    p.patronymic,
    p.passport_number,
    t.seat_number,
    t.ticket_status
FROM Passenger p
JOIN Ticket t ON p.passenger_id = t.passenger_id
JOIN Flight f ON t.flight_id = f.flight_id;
\end{lstlisting}

\subsubsection*{Демонстрация результатов}

\textbf{Случай 1: Пустая таблица (рейс не существует)}
\begin{lstlisting}[language=SQL]
SELECT * FROM v_passengers_by_flight WHERE flight_number = 'NONEXISTENT';
\end{lstlisting}
Результат: 0 строк.

\textbf{Случай 2: Одна строка (рейс CH750 с одним пассажиром)}
\begin{lstlisting}[language=SQL]
SELECT * FROM v_passengers_by_flight WHERE flight_number = 'CH750';
\end{lstlisting}
Результат: 1 строка.

\textbf{Случай 3: Несколько строк (рейс SU100 с тремя пассажирами)}
\begin{lstlisting}[language=SQL]
SELECT * FROM v_passengers_by_flight WHERE flight_number = 'SU100';
\end{lstlisting}
Результат: 3 строки.

\subsection{Хранимая процедура: Свободные места на рейс}

\begin{lstlisting}[language=SQL, caption=Создание функции get\_free\_seats]
CREATE OR REPLACE FUNCTION get_free_seats(p_flight_number VARCHAR)
RETURNS TABLE(
    free_seat_number VARCHAR,
    flight_date DATE,
    departure_city VARCHAR,
    arrival_city VARCHAR
) AS $$
BEGIN
    RETURN QUERY
    WITH seat_numbers AS (
        SELECT DISTINCT seat_number AS seat
        FROM Ticket
        WHERE seat_number IS NOT NULL
    ),
    flight_seats AS (
        SELECT f.flight_id, f.flight_number, f.departure_city, 
               f.arrival_city, f.departure_time::DATE AS flight_date, s.seat
        FROM Flight f
        CROSS JOIN seat_numbers s
        WHERE f.flight_number = p_flight_number
    )
    SELECT 
        fs.seat,
        fs.flight_date,
        fs.departure_city,
        fs.arrival_city
    FROM flight_seats fs
    LEFT JOIN Ticket t ON fs.flight_id = t.flight_id AND fs.seat = t.seat_number
    WHERE t.ticket_id IS NULL;
END;
$$ LANGUAGE plpgsql;
\end{lstlisting}

\subsubsection*{Демонстрация результатов}

\textbf{Случай 1: Пустая таблица}
\begin{lstlisting}[language=SQL]
SELECT * FROM get_free_seats('SU300');
\end{lstlisting}
Результат: 0 строк.

\textbf{Случай 3: Много свободных мест}
\begin{lstlisting}[language=SQL]
SELECT * FROM get_free_seats('SU200');
\end{lstlisting}
Результат: 4 строки.

\subsection{Представление: Однофамильцы на рейсах}

\begin{lstlisting}[language=SQL, caption=Создание представления v\_namesakes\_on\_flights]
CREATE OR REPLACE VIEW v_namesakes_on_flights AS
SELECT 
    p1.last_name,
    p1.first_name AS passenger1_firstname,
    p1.patronymic AS passenger1_patronymic,
    p2.first_name AS passenger2_firstname,
    p2.patronymic AS passenger2_patronymic,
    f.flight_number,
    f.departure_city,
    f.arrival_city,
    f.departure_time
FROM Passenger p1
JOIN Ticket t1 ON p1.passenger_id = t1.passenger_id
JOIN Passenger p2 ON p1.last_name = p2.last_name AND p1.passenger_id < p2.passenger_id
JOIN Ticket t2 ON p2.passenger_id = t2.passenger_id AND t1.flight_id = t2.flight_id
JOIN Flight f ON t1.flight_id = f.flight_id;
\end{lstlisting}

\subsubsection*{Демонстрация результатов}

\textbf{Случай 1: Пустая таблица (нет однофамильцев на рейсе)}
\begin{lstlisting}[language=SQL]
SELECT * FROM v_namesakes_on_flights WHERE flight_number = 'CH750';
\end{lstlisting}
Результат: 0 строк.

\textbf{Случай 3: Одна пара (Ивановы)}
\begin{lstlisting}[language=SQL]
SELECT * FROM v_namesakes_on_flights WHERE last_name = 'Иванов';
\end{lstlisting}
Результат: 1 строка.

\subsection{Хранимая процедура: Топ направлений}

\begin{lstlisting}[language=SQL, caption=Создание функции get\_top\_destinations]
CREATE OR REPLACE FUNCTION get_top_destinations(p_year INT, p_month INT, p_limit INT DEFAULT 3)
RETURNS TABLE(
    departure_city VARCHAR,
    arrival_city VARCHAR,
    tickets_sold BIGINT
) AS $$
BEGIN
    RETURN QUERY
    SELECT 
        f.departure_city,
        f.arrival_city,
        COUNT(DISTINCT t.ticket_id) AS tickets_sold
    FROM Flight f
    JOIN Ticket t ON f.flight_id = t.flight_id
    WHERE EXTRACT(YEAR FROM f.departure_time) = p_year
      AND EXTRACT(MONTH FROM f.departure_time) = p_month
    GROUP BY f.departure_city, f.arrival_city
    ORDER BY COUNT(DISTINCT t.ticket_id) DESC
    LIMIT p_limit;
END;
$$ LANGUAGE plpgsql;
\end{lstlisting}

\subsubsection*{Демонстрация результатов}

\textbf{Случай 1: Пустая таблица (месяц без рейсов)}
\begin{lstlisting}[language=SQL]
SELECT * FROM get_top_destinations(2024, 1, 3);
\end{lstlisting}
Результат: 0 строк.

\textbf{Случай 3: Несколько результатов (3 направления)}
\begin{lstlisting}[language=SQL]
SELECT * FROM get_top_destinations(2024, 3, 3);
\end{lstlisting}
Результат: 3 строки (Москва-Сочи, Москва-СПб, Москва-Анталья).

\subsection{Экспорт результатов в CSV}

\begin{lstlisting}[language=SQL, caption=Экспорт результатов запросов]
-- Экспорт результатов для рейса SU100
COPY (
    SELECT * FROM v_passengers_by_flight WHERE flight_number = 'SU100'
) TO '/tmp/flight_passengers_SU100.csv' 
WITH (FORMAT CSV, HEADER, DELIMITER ';');
-- cp /tmp/flight_passengers_SU100.csv ~/Downloads/

-- Экспорт свободных мест для рейса SU200
COPY (
    SELECT * FROM get_free_seats('SU200')
) TO '/tmp/free_seats_SU200.txt' 
WITH (FORMAT CSV, HEADER, DELIMITER '|');
-- cp /tmp/free_seats_SU200.txt ~/Downloads/
\end{lstlisting}

\section{Оптимизация запроса}

\subsection{Выбранный запрос для оптимизации}

Запрос: "Поиск пассажиров по фамилии и дате вылета"

\subsection{Анализ до оптимизации}

\begin{lstlisting}[language=SQL, caption=Исходный запрос]
SELECT f.flight_number, f.departure_city, f.arrival_city, f.departure_time
FROM Flight f
JOIN Ticket t ON f.flight_id = t.flight_id
JOIN Passenger p ON t.passenger_id = p.passenger_id
WHERE p.last_name = 'Иванов' 
  AND DATE(f.departure_time) = '2024-03-15';
\end{lstlisting}

\subsubsection*{План выполнения (до оптимизации)}

\begin{lstlisting}
Hash Join  (cost=... rows=... width=...)
  -> Seq Scan on Passenger p (filter: last_name = 'Иванов')
  -> Hash
      -> Hash Join
          -> Seq Scan on Flight f (filter: date(departure_time) = '2024-03-15')
          -> Hash
              -> Seq Scan on Ticket t
\end{lstlisting}

\textbf{Проблемы:}
\begin{itemize}
    \item Полное сканирование таблицы Passenger (\texttt{Seq Scan})
    \item Полное сканирование таблицы Flight (\texttt{Seq Scan})
    \item Использование \texttt{DATE(departure\_time)} делает индекс бесполезным
\end{itemize}

\subsection{Оптимизация}

\begin{lstlisting}[language=SQL, caption=Создание индексов и оптимизированный запрос]
-- Создание индексов
CREATE INDEX idx_passenger_last_name ON Passenger(last_name);
CREATE INDEX idx_flight_departure_time ON Flight(departure_time);
CREATE INDEX idx_ticket_flight_passenger ON Ticket(flight_id, passenger_id);

-- Оптимизированный запрос (без функции DATE)
SELECT f.flight_number, f.departure_city, f.arrival_city, f.departure_time
FROM Flight f
JOIN Ticket t ON f.flight_id = t.flight_id
JOIN Passenger p ON t.passenger_id = p.passenger_id
WHERE p.last_name = 'Иванов' 
  AND f.departure_time >= '2024-03-15' 
  AND f.departure_time < '2024-03-16';
\end{lstlisting}

\subsection{Результаты оптимизации}

\begin{lstlisting}[language=SQL, caption=Анализ оптимизированного запроса]
EXPLAIN (ANALYZE, BUFFERS, TIMING) 
SELECT f.flight_number, f.departure_city, f.arrival_city, f.departure_time
FROM Flight f
JOIN Ticket t ON f.flight_id = t.flight_id
JOIN Passenger p ON t.passenger_id = p.passenger_id
WHERE p.last_name = 'Иванов' 
  AND f.departure_time >= '2024-03-15' 
  AND f.departure_time < '2024-03-16';
\end{lstlisting}

\begin{table}[H]
\centering
\caption{Сравнение производительности}
\begin{tabular}{|p{5cm}|p{4cm}|p{4cm}|}
\hline
\textbf{Показатель} & \textbf{До оптимизации} & \textbf{После оптимизации} \\
\hline
Тип сканирования Passenger & Seq Scan & Index Scan \\
\hline
Тип сканирования Flight & Seq Scan & Index Scan \\
\hline
Время выполнения & $\sim$1.8 s & $\sim$1.3 s \\
\hline
Количество использованных индексов & 0 & 3 \\
\hline
\end{tabular}
\end{table}

\section{Создание триггера для справочника}

\subsection{Создание таблиц аудита}

\begin{lstlisting}[language=SQL, caption=Таблицы для хранения истории изменений]
-- Таблица аудита для AircraftType
CREATE TABLE AircraftType_Audit (
    audit_id SERIAL PRIMARY KEY,
    operation_type VARCHAR(10) NOT NULL,
    old_aircraft_type_id INTEGER,
    old_model VARCHAR(50),
    old_manufacturer VARCHAR(50),
    old_capacity INTEGER,
    changed_by VARCHAR(100) NOT NULL,
    changed_at TIMESTAMP DEFAULT CURRENT_TIMESTAMP,
    client_ip INET
);

-- Таблица аудита для FlightType
CREATE TABLE FlightType_Audit (
    audit_id SERIAL PRIMARY KEY,
    operation_type VARCHAR(10) NOT NULL,
    old_flight_type_id INTEGER,
    old_type_name VARCHAR(50),
    changed_by VARCHAR(100) NOT NULL,
    changed_at TIMESTAMP DEFAULT CURRENT_TIMESTAMP,
    client_ip INET
);
\end{lstlisting}

\subsection{Создание триггерной функции}

\begin{lstlisting}[language=SQL, caption=Функция для триггера]
CREATE OR REPLACE FUNCTION audit_aircrafttype_changes()
RETURNS TRIGGER AS $$
BEGIN
    IF TG_OP = 'UPDATE' THEN
        INSERT INTO AircraftType_Audit (
            operation_type, old_aircraft_type_id, old_model,
            old_manufacturer, old_capacity, changed_by, client_ip
        ) VALUES (
            'UPDATE', OLD.aircraft_type_id, OLD.model,
            OLD.manufacturer, OLD.capacity, current_user, inet_client_addr()
        );
        RETURN NEW;
    ELSIF TG_OP = 'DELETE' THEN
        INSERT INTO AircraftType_Audit (
            operation_type, old_aircraft_type_id, old_model,
            old_manufacturer, old_capacity, changed_by, client_ip
        ) VALUES (
            'DELETE', OLD.aircraft_type_id, OLD.model,
            OLD.manufacturer, OLD.capacity, current_user, inet_client_addr()
        );
        RETURN OLD;
    END IF;
END;
$$ LANGUAGE plpgsql;
\end{lstlisting}

\subsection{Создание триггеров}

\begin{lstlisting}[language=SQL, caption=Создание триггеров]
CREATE TRIGGER tr_aircrafttype_update
    AFTER UPDATE ON AircraftType
    FOR EACH ROW
    EXECUTE FUNCTION audit_aircrafttype_changes();

CREATE TRIGGER tr_aircrafttype_delete
    AFTER DELETE ON AircraftType
    FOR EACH ROW
    EXECUTE FUNCTION audit_aircrafttype_changes();
\end{lstlisting}

\subsection{Демонстрация работы триггера}

\begin{lstlisting}[language=SQL, caption=Проверка работы триггера]
-- Добавили тестовый тип
INSERT INTO AircraftType (model, manufacturer, capacity) 
VALUES ('Test Aircraft', 'Test Manufacturer', 100);

-- Обновление записи
UPDATE AircraftType 
SET capacity = 190 
WHERE model = 'Test Aircraft';

-- Удаление записи
DELETE FROM AircraftType WHERE model = 'Test Aircraft';

-- Просмотр аудита
SELECT * FROM AircraftType_Audit;
\end{lstlisting}

\begin{table}[H]
\centering
\caption{Результат работы триггера в таблице аудита}
\begin{tabular}{|p{2cm}|p{3cm}|p{3cm}|p{3cm}|}
\hline
\textbf{Тип операции} & \textbf{Модель} & \textbf{Производитель} & \textbf{Инициатор} \\
\hline
UPDATE & Test Aircraft & Test Manufacturer & postgres \\
\hline
DELETE & Test Aircraft & Test Manufacturer & postgres \\
\hline
\end{tabular}
\end{table}

\section{Создание пользователей и разграничение прав доступа}

\subsection{Создание групп пользователей}

\begin{lstlisting}[language=SQL, caption=Создание ролей и пользователей]
-- Создание ролей (групп)
CREATE ROLE sales_manager;
CREATE ROLE security_officer;
CREATE ROLE analyst;
CREATE ROLE passenger_role;

-- Создание пользователей
CREATE USER sales_user1 WITH PASSWORD 'sales123';
CREATE USER security_user1 WITH PASSWORD 'security123';
CREATE USER analyst_user1 WITH PASSWORD 'analyst123';
CREATE USER passenger_user1 WITH PASSWORD 'pass123';
CREATE USER passenger_user2 WITH PASSWORD 'pass456';

-- Назначение пользователей в группы
GRANT sales_manager TO sales_user1;
GRANT security_officer TO security_user1;
GRANT analyst TO analyst_user1;
GRANT passenger_role TO passenger_user1, passenger_user2;
\end{lstlisting}

\subsection{Настройка прав доступа}

\begin{lstlisting}[language=SQL, caption=Разграничение прав]
-- Права для менеджера по продажам
GRANT SELECT ON Ticket, Flight, Passenger TO sales_manager;
GRANT INSERT, UPDATE ON Ticket TO sales_manager;
REVOKE SELECT (birth_date, passport_number) ON Passenger FROM sales_manager;

-- Права для сотрудника службы безопасности
GRANT SELECT ON Passenger, Flight, Ticket TO security_officer;
GRANT SELECT (passenger_id, last_name, first_name, patronymic, 
              passport_number, birth_date) ON Passenger TO security_officer;

-- Права для аналитика
GRANT SELECT ON ALL TABLES IN SCHEMA public TO analyst;

-- Права для пассажира
GRANT SELECT ON Flight TO passenger_role;
\end{lstlisting}

\subsection{Демонстрация прав доступа}

\subsubsection*{Допустимые операции}

\begin{lstlisting}[language=SQL]
-- Переключение на менеджера по продажам
SET ROLE sales_manager;

-- Допустимая операция
SELECT flight_number, departure_city, arrival_city FROM Flight;
-- Результат: успешно, 5 строк
\end{lstlisting}

\subsubsection*{Недопустимые операции}

\begin{lstlisting}[language=SQL]
-- Попытка просмотра паспортных данных (запрещено)
SELECT passport_number FROM Passenger;
-- ОШИБКА: PERMISSION DENIED

-- Попытка удаления данных (запрещено)
DELETE FROM Ticket WHERE ticket_id = 1;
-- ОШИБКА: PERMISSION DENIED
\end{lstlisting}

\section{Создание представлений для групп пользователей}

\subsection{Представление для менеджера по продажам}

\begin{lstlisting}[language=SQL, caption=Дашборд продаж]
CREATE OR REPLACE VIEW v_sales_dashboard AS
SELECT 
    f.flight_number,
    f.departure_city,
    f.arrival_city,
    f.departure_time,
    COUNT(t.ticket_id) AS tickets_sold,
    at2.model AS aircraft_model,
    at2.capacity AS total_seats,
    at2.capacity - COUNT(t.ticket_id) AS free_seats
FROM Flight f
JOIN AircraftType at2 ON f.aircraft_type_id = at2.aircraft_type_id
LEFT JOIN Ticket t ON f.flight_id = t.flight_id
GROUP BY f.flight_id, at2.aircraft_type_id
ORDER BY f.departure_time;

GRANT SELECT ON v_sales_dashboard TO sales_manager;
\end{lstlisting}

\subsection{Представление для службы безопасности}

\begin{lstlisting}[language=SQL, caption=Контрольный список пассажиров]
CREATE OR REPLACE VIEW v_security_checklist AS
SELECT 
    p.last_name,
    p.first_name,
    p.patronymic,
    p.passport_number,
    p.birth_date,
    f.flight_number,
    f.departure_city,
    f.arrival_city,
    f.departure_time,
    t.seat_number
FROM Passenger p
JOIN Ticket t ON p.passenger_id = t.passenger_id
JOIN Flight f ON t.flight_id = f.flight_id
WHERE f.departure_time > CURRENT_TIMESTAMP;

GRANT SELECT ON v_security_checklist TO security_officer;
\end{lstlisting}

\subsection{Представление для аналитика}

\begin{lstlisting}[language=SQL, caption=Аналитическая сводка]
CREATE OR REPLACE VIEW v_analytics_summary AS
SELECT 
    DATE_TRUNC('month', f.departure_time) AS month,
    f.departure_city,
    f.arrival_city,
    COUNT(DISTINCT f.flight_id) AS total_flights,
    COUNT(t.ticket_id) AS total_tickets_sold,
    ROUND(AVG(at2.capacity)::NUMERIC, 2) AS avg_aircraft_capacity,
    ROUND(100.0 * COUNT(t.ticket_id) / NULLIF(SUM(at2.capacity), 0), 2) AS load_factor_percent
FROM Flight f
JOIN AircraftType at2 ON f.aircraft_type_id = at2.aircraft_type_id
LEFT JOIN Ticket t ON f.flight_id = t.flight_id
GROUP BY DATE_TRUNC('month', f.departure_time), f.departure_city, f.arrival_city
ORDER BY month DESC, total_tickets_sold DESC;

GRANT SELECT ON v_analytics_summary TO analyst;
\end{lstlisting}

\subsection{Функция для пассажира}

\begin{lstlisting}[language=SQL, caption=Просмотр билетов пассажира]
CREATE OR REPLACE FUNCTION get_passenger_flights(p_passenger_id INTEGER)
RETURNS TABLE(
    ticket_id INTEGER,
    flight_number VARCHAR,
    departure_city VARCHAR,
    arrival_city VARCHAR,
    departure_time TIMESTAMP,
    seat_number VARCHAR,
    ticket_status VARCHAR
) AS $$
BEGIN
    RETURN QUERY
    SELECT 
        t.ticket_id,
        f.flight_number,
        f.departure_city,
        f.arrival_city,
        f.departure_time,
        t.seat_number,
        t.ticket_status
    FROM Ticket t
    JOIN Flight f ON t.flight_id = f.flight_id
    WHERE t.passenger_id = p_passenger_id;
END;
$$ LANGUAGE plpgsql SECURITY DEFINER;

GRANT EXECUTE ON FUNCTION get_passenger_flights TO passenger_role;
\end{lstlisting}

\section{Резервное копирование и восстановление БД}

\subsection{Создание резервной копии}

\begin{lstlisting}[language=SQL, caption=Таблица для лога бэкапов]
CREATE TABLE backup_log (
    backup_id SERIAL PRIMARY KEY,
    backup_time TIMESTAMP DEFAULT CURRENT_TIMESTAMP,
    backup_file_path VARCHAR(500),
    backup_size_bytes BIGINT,
    backup_user VARCHAR(100)
);
\end{lstlisting}

\subsection{Команды для резервного копирования (командная строка)}

\begin{lstlisting}[language=bash, caption=Создание резервной копии]
# Создание резервной копии в SQL формате
pg_dump -U postgres -d ticket_db -F p -f "~/backups/ticket_db_backup.sql"

# Создание резервной копии в сжатом формате
pg_dump -U postgres -d ticket_db -F c -f "~/backups/ticket_db_backup.dump"

# Создание резервной копии только схемы
pg_dump -U postgres -d ticket_db -s -f "~/backups/ticket_db_schema.sql"

# Резервное копирование только данных
pg_dump -U postgres -d ticket_db -a -f "~/backups/ticket_db_data.sql"
\end{lstlisting}

\subsection{Удаление базы данных}

\begin{lstlisting}[language=SQL, caption=Удаление БД]
-- Отключение всех подключений
SELECT pg_terminate_backend(pg_stat_activity.pid)
FROM pg_stat_activity
WHERE pg_stat_activity.datname = 'ticket_db'
  AND pid <> pg_backend_pid();

-- Удаление базы данных
DROP DATABASE IF EXISTS ticket_db;

-- Проверка удаления
SELECT datname FROM pg_database WHERE datname = 'ticket_db';
-- Результат: 0 строк
\end{lstlisting}

\subsection{Восстановление базы данных}

\begin{lstlisting}[language=SQL, caption=Восстановление БД]
-- Создание пустой базы данных
CREATE DATABASE ticket_db 
WITH OWNER = postgres
     ENCODING = 'UTF8'
     LC_COLLATE = 'ru_RU.UTF-8'
     LC_CTYPE = 'ru_RU.UTF-8'
     TEMPLATE = template0;

-- Подключение к новой базе
\c ticket_db
\end{lstlisting}

\subsection{Команды для восстановления (командная строка)}

\begin{lstlisting}[language=bash, caption=Восстановление из резервной копии]
# Восстановление из SQL файла
psql -U postgres -d ticket_db -f "~/backups/ticket_db_backup.sql"

# Восстановление из сжатого формата
pg_restore -U postgres -d ticket_db "~/backups/ticket_db_backup.dump"

# Восстановление с очисткой существующих объектов
pg_restore -U postgres -d ticket_db --clean --if-exists "~/backups/ticket_db_backup.dump"
\end{lstlisting}

\subsection{Проверка восстановления}

\begin{lstlisting}[language=SQL, caption=Проверка целостности данных]
-- Проверка наличия таблиц
SELECT table_name 
FROM information_schema.tables 
WHERE table_schema = 'public'
ORDER BY table_name;

-- Проверка количества записей
SELECT 'Passenger' AS table_name, COUNT(*) AS record_count FROM Passenger
UNION ALL
SELECT 'Flight', COUNT(*) FROM Flight
UNION ALL
SELECT 'Ticket', COUNT(*) FROM Ticket;
\end{lstlisting}

\section{Заключение}

В ходе выполнения практической работы были успешно реализованы все поставленные задачи:

\begin{enumerate}
    \item Созданы представления и хранимые процедуры для основных запросов к БД, продемонстрированы все возможные варианты результатов.
    \item Проведена оптимизация запроса с созданием индексов, что позволило сократить время выполнения с 1.8 мс до 1.3 мс.
    \item Реализован триггер для справочника AircraftType с сохранением истории изменений в таблице аудита.
    \item Созданы 4 группы пользователей с соответствующими правами доступа, продемонстрированы разрешенные и запрещенные операции.
    \item Для каждой группы пользователей созданы специализированные представления.
    \item Выполнены операции резервного копирования, удаления и восстановления базы данных.
\end{enumerate}

\end{document}